\chapter{移位寄存器考试变式题}

\section*{题目1：串并转换（SIPO）}

\textbf{题目}：设计一个4位串入并出移位寄存器。串行输入数据，4个时钟周期后并行输出完整4位数据。

\textbf{VHDL}：
\begin{lstlisting}
signal reg : std_logic_vector(3 downto 0);
process(clk) begin
  if rising_edge(clk) then
    reg <= si & reg(3 downto 1);  -- 串行右移
  end if;
end process;
po <= reg;
\end{lstlisting}

\textbf{逐拍分析（输入1011）}：
\begin{center}
\begin{tabular}{|c|c|c|}
\hline
拍 & 输入 & [Q3 Q2 Q1 Q0] \\
\hline
1 & 1 & [1 0 0 0] \\
2 & 0 & [0 1 0 0] \\
3 & 1 & [1 0 1 0] \\
4 & 1 & [1 1 0 1] \\
\hline
\end{tabular}
\end{center}

4拍后po=1101（数据反转了，因为右移）。如想保持顺序，用左移。

\section*{题目2：并串转换（PISO）}

\textbf{题目}：设计4位并入串出移位寄存器。先load并行数据，然后移位输出。

\textbf{VHDL}：
\begin{lstlisting}
signal reg : std_logic_vector(3 downto 0);
process(clk) begin
  if rising_edge(clk) then
    if load = '1' then
      reg <= data_in;
    else
      reg <= '0' & reg(3 downto 1);  -- 右移，高位补0
    end if;
  end if;
end process;
so <= reg(0);
\end{lstlisting}

\section*{题目3：双向移位寄存器}

\textbf{题目}：设计4位双向移位寄存器。dir=1右移，dir=0左移。

\begin{lstlisting}
 if dir = '1' then
  reg <= si_right & reg(3 downto 1);
else
  reg <= reg(2 downto 0) & si_left;
end if;
\end{lstlisting}

\section*{题目4：带并行装载的通用移位寄存器（74LS194风格）}

\textbf{题目}：设计类似74LS194的4位通用移位寄存器。支持：保持、右移、左移、并行装载四种模式。

\textbf{VHDL}：
\begin{lstlisting}
case mode is
  when "00" => null;                    -- 保持
  when "01" => reg <= si_r & reg(3 downto 1);  -- 右移
  when "10" => reg <= reg(2 downto 0) & si_l;  -- 左移
  when "11" => reg <= data_in;          -- 装载
end case;
\end{lstlisting}

\section*{题目5：桶形移位器}

\textbf{题目}：设计8位桶形移位器，一次可移0-7位（移位量由3位输入控制）。

\textbf{分析}：桶形移位器是纯组合逻辑（不需要时钟），本质上是多路选择器网络。

\section*{题目6：移位寄存器的延迟线应用}

\textbf{题目}：用移位寄存器实现4拍延迟线（输入信号延迟4个时钟周期后输出）。

\textbf{解}：4位移位寄存器，so输出 = 4拍前的si输入。

\section*{题目7：序列检测用移位寄存器方案}

\textbf{题目}：用移位寄存器+比较器检测序列1011（替代状态机方案）。

\textbf{分析}：4位移位寄存器存储最近4位输入。当reg="1011"时输出1。

\begin{lstlisting}
reg <= data_in & reg(3 downto 1);
detected <= '1' when reg = "1011" else '0';
\end{lstlisting}

\section*{题目8：8位移位寄存器}

\textbf{题目}：设计8位串入并出移位寄存器。

\textbf{答案}：位宽从4扩展到8，其他逻辑完全相同。

\section*{题目9：循环移位寄存器}

\textbf{题目}：设计循环移位寄存器（最低位移入最高位，不丢失数据）。

\textbf{VHDL}：
\begin{lstlisting}
reg <= reg(0) & reg(7 downto 1);  -- 循环右移
\end{lstlisting}

\section*{题目10：算术移位寄存器}

\textbf{题目}：设计算术右移寄存器（保持符号位不变）。

\begin{lstlisting}
reg <= reg(7) & reg(7 downto 1);  -- 算术右移：符号位不变
\end{lstlisting}

\section*{题目11：逻辑移位vs算术移位}

\textbf{题目}：比较逻辑左移和算术左移的区别。

\textbf{分析}：左移时逻辑移位和算术移位等价（都补0）。右移时：逻辑右移补0，算术右移补符号位。

\section*{题目12：移位寄存器实现乘除法}

\textbf{题目}：用移位寄存器实现×2和÷2运算。

\textbf{分析}：
\begin{itemize}
  \item 左移1位 = ×2（reg <= reg(6 downto 0) \& '0'）
  \item 右移1位 = ÷2（reg <= '0' \& reg(7 downto 1)）
\end{itemize}

\section*{题目13：给定移位寄存器初值和输入序列，推断最终状态}

\textbf{题目}：4位右移寄存器初值0000，输入序列1011。4拍后寄存器内容？

\textbf{解（逐拍）}：
\begin{itemize}
  \item 初始：[0000]
  \item 拍1(si=1)：[1000]
  \item 拍2(si=0)：[0100]
  \item 拍3(si=1)：[1010]
  \item 拍4(si=1)：[1101]
\end{itemize}

4拍后：[1101]

\section*{题目14：移位寄存器Testbench设计}

\textbf{题目}：为8位移位寄存器设计完整的Testbench，验证串入并出功能。

\section*{题目15：移位寄存器构建伪随机序列（LFSR）}

\textbf{题目}：用4位移位寄存器+反馈构建最大长度LFSR。

\textbf{分析}：反馈多项式$X^4 + X^3 + 1$。反馈=Q3 XOR Q2。
产生序列长度=$2^4-1=15$（全0状态除外）。

\begin{examtip}[移位寄存器类题目的统一解法]
\begin{enumerate}
  \item \textbf{右移：}\texttt{reg <= si \& reg(N-1 downto 1)}
  \item \textbf{左移：}\texttt{reg <= reg(N-2 downto 0) \& si}
  \item \textbf{循环：}\texttt{reg <= reg(0) \& reg(N-1 downto 1)}
  \item \textbf{并行装载：}\texttt{if load='1' then reg <= data\_in;}
  \item \textbf{串出：}reg(0)（右移）或 reg(N-1)（左移）
  \item \textbf{并出：}整个reg
\end{enumerate}
\end{examtip}
